\section{The arithmetic question behind the experiments}
For $s>1$, Dirichlet's beta function is
\[
 \beta(s)=\sum_{j=0}^{\infty}\frac{(-1)^j}{(2j+1)^s}.
\]
Its positive odd values are nonzero rational multiples of powers of
$\pi$. Individual even values are much harder to study. Write
$G=\beta(2)$ for Catalan's constant. This supplement concerns attempts
to prove irrationality of $\beta(4)$, of a member of a very small set
of even beta values, or of combinations $G+\lambda\pi^2$ with rational
$\lambda$.

The earlier investigation supplement~\cite{firstsupp} treats the first
collection of determinant and integral approaches. The present document
covers all 29 subsequent research records from 6--8 November 2026.
The finite-count and
shifted-interval theorems~\cite{count,shifted} provide context, but their
proofs are not repeated here. The question now is whether their techniques,
or other rational approximations, can reach a specified individual value
or a much smaller set.

The intended reader is familiar with partial fractions, prime valuations,
basic complex analysis and finite-dimensional linear algebra. The Gram
and lattice arguments are introduced through the particular quantity
they must control. The first two constructions can be read without the
longer mixed-moment analysis.

\subsection{What would prove irrationality?}
Suppose a fixed finite set of real numbers $\alpha_1,\ldots,\alpha_k$
were rational. A common denominator $D$ would make
\[
 D\bigl(a_0+a_1\alpha_1+\cdots+a_k\alpha_k\bigr)\in\Z
 \qquad(a_0,\ldots,a_k\in\Z).
\]
Thus nonzero integer-coefficient forms tending to zero would contradict
that assumption. This elementary observation puts three requirements
on a construction: the right support, enough smallness after the
coefficients are made integral, and nonvanishing along an unbounded
sequence.

A small approximation error can conceal the second requirement.
For example, $|\beta(4)-p/q|$ must be multiplied by $q$ before it becomes
the value of the integer form $q\beta(4)-p$. A spectacularly accurate
fraction can still give a large integer form. Several experiments below
improve the approximation while falling short at precisely this step.

\begin{definition}
For a nonzero rational vector $v$, its rational content $\cont(v)>0$
is the positive rational number for which $v/\cont(v)$ is a primitive
integer vector. For a rational polynomial, apply this definition to
its coefficient vector. Its primitive normalization divides by this
content, with no change of sign.
\end{definition}

\begin{lemma}[Separate the rational constant's cost]\label{lem:constant-cost}
Let $q+\sum_j\rho_j\alpha_j$ have some $\rho_j\ne0$. Put
$g=\cont((\rho_j)_j)$, $A_j=\rho_j/g$, and write $q/g=a/d$ in lowest
terms with $d>0$. The full primitive vector is
\[
             (a,dA_1,\ldots,dA_k),
\]
and its multiplier relative to the original form is $d/g$.
\end{lemma}
\begin{proof}
The integer vector $(A_j)$ has gcd one, so $(dA_j)$ has gcd $d$.
Since $\gcd(a,d)=1$, adjoining $a$ leaves gcd one.
\end{proof}

The lemma will explain a concrete failure in Section~\ref{sec:small}:
two expressions have identical beta coefficients but different rational
constants. Removing a finite prefix improves the real size, yet the new
constant has a larger denominator after normalization by $g$.

For determinant polynomials the same principle applies to every
coefficient, rather than just a pair. If $P_n\in\Z[X]$, $\deg P_n=O(n)$,
and $\theta=a/d$ is a fixed rational number, then
$d^{\deg P_n}P_n(\theta)$ is integral. On a logarithmic $n^2$ scale,
this fixed-denominator cost is $O(n)$ and is negligible. Allowing $d$
to grow with $n$ destroys that argument. This distinction matters in
the rational test evaluations of Section~\ref{sec:corrections}.

\subsection{How to read the results}
The main lines of investigation are summarized below. A limitation of
one specified estimate is not a limitation of every estimate for the
same construction.
\begin{center}
\begin{tabular}{@{}>{\raggedright\arraybackslash}p{.22\textwidth}>{\raggedright\arraybackslash}p{.37\textwidth}>{\raggedright\arraybackslash}p{.34\textwidth}@{}}
\toprule
Construction & Established result & Remaining issue \\
\midrule
Odd-grid isolation & Fixed polynomial multipliers cannot isolate
$\beta(4)$ under the stated degree condition. & Other multipliers,
families or nonvanishing mechanisms. \\[3pt]
Real Pad\'e forms & Exact real remainders and their exponential rates. &
Further common coefficient factors or different evaluation regimes. \\[3pt]
Mixed moments & Proved decay is sufficient before integer corrections. &
Integer corrections that preserve enough of that decay. \\[3pt]
Small beta sets & Complete forms and an exact large-prime residue formula. &
A useful arithmetic balance beyond the finite tests. \\
\bottomrule
\end{tabular}
\end{center}

Sections~\ref{sec:isolation} and~\ref{sec:pade} can be read independently.
Sections~\ref{sec:moments}--\ref{sec:corrections} develop the mixed-moment
problem in order: form a matrix whose entries are linear in $G+\pi^2$,
bound its determinant, then identify what the integer corrections
must accomplish.
Section~\ref{sec:small} returns to the two- and three-value targets.
Section~\ref{sec:variants} treats other tested variants. The catalog in
Appendix~\ref{app:catalog} records the shorter investigations and
superseded intermediate bounds, with an outcome for each source record.

\subsection{Sources, certificates and review scope}
The odd-grid construction is inherited from~\cite{shifted}, with the
rational-function lineage described there. The real Pad\'e construction
comes from Kawashima--Po\"els~\cite{KP}; its real interpretation follows
the Cauchy-transform viewpoint of Rivoal~\cite{Rivoal06}. The mixed
moments and Chebyshev rows come from OpenAI~\cite{OpenAI}. The small-set
derivative construction follows the approach of Krisman~\cite{Krisman},
including his use of triple zeros for grouped denominator cancellation.
Precise attributions are repeated at the relevant proof steps.

The accompanying files give the selected calculations and a complete
record catalog. Classical inputs and imported construction lemmas are
identified in the text. Rational and interval computations have distinct
roles: an exact coefficient identity, a certified finite inequality and
a limiting theorem are separate assertions.

The underlying canonical decay estimate and this supplement received
separate independent mathematical audits. Appendix~\ref{app:audit} gives
their scopes, records the missing hypothesis identified in the correction
argument, and describes the accompanying checks.
